% Original English questions and concise English solutions. Compile using XeLaTeX.
\documentclass[UTF8,10pt,letterpaper,fontset=fandol]{ctexart}
\usepackage[margin=1in]{geometry}
\usepackage{amsmath,amssymb}
\usepackage{needspace}
\usepackage[unicode,hidelinks]{hyperref}
\pagestyle{plain}
\raggedbottom
\setlength{\parindent}{0pt}
\setlength{\parskip}{1pt}
\setlength{\emergencystretch}{2em}
\setcounter{secnumdepth}{0}
\ctexset{section={format=\Large\bfseries,beforeskip=1.2ex,afterskip=0.8ex}}
\newcommand{\F}{\mathbb F}
\newcommand{\Solution}{\par\Needspace{4\baselineskip}\textbf{Solution}\par}
\begin{document}
\begin{center}
{\LARGE Homework 4}\\[0.6em]
October 16, 2026\\[0.25em]
Name: YuanHao Tang\quad Student ID: 3120265621
\end{center}

\section{Problem 1}
Find out the Smith canonical form of the $\lambda$-matrix
\[
A(\lambda)=\begin{pmatrix}
1-\lambda&\lambda^2&\lambda\\
2\lambda-1&\lambda(1-\lambda)&-2\lambda\\
1+\lambda^2&\lambda^2&-\lambda^2
\end{pmatrix}.
\]
\Solution
Let $D_k$ be the monic gcd of all $k\times k$ minors, with $D_0=1$. Then the invariant polynomials are $d_k=D_k/D_{k-1}$.
Since $2(1-\lambda)+(2\lambda-1)=1$, $D_1=1$. Columns two and three are divisible by $\lambda$, so every second-order minor is divisible by $\lambda$. Moreover,
\[
\begin{vmatrix}1-\lambda&\lambda\\2\lambda-1&-2\lambda\end{vmatrix}=-\lambda,
\]
giving $D_2=\lambda$. Since $\det A=-\lambda^2(\lambda+1)$, $D_3=\lambda^2(\lambda+1)$. Thus
\[
\boxed{S_A(\lambda)=\operatorname{diag}\bigl(1,\lambda,\lambda(\lambda+1)\bigr)}.
\]

\section{Problem 2}
Let $A(\lambda)$ be the $n\times n$ $\lambda$-matrix
\[
A(\lambda)=\begin{pmatrix}
\lambda-a&b_1&0&\cdots&0\\
0&\lambda-a&b_2&\ddots&\vdots\\
\vdots&\ddots&\ddots&\ddots&0\\
0&\cdots&0&\lambda-a&b_{n-1}\\
0&\cdots&\cdots&0&\lambda-a
\end{pmatrix},
\]
where $b_1,b_2,\ldots,b_{n-1}$ are all nonzero constants. Find out the invariant polynomials of $A(\lambda)$.
\Solution
For $1\le k<n$, the minor on rows $1,\ldots,k$ and columns $2,\ldots,k+1$ is the nonzero constant $b_1\cdots b_k$, so $D_k=1$.
Also $D_n=\det A=(\lambda-a)^n$. Therefore
\[
\boxed{d_1=\cdots=d_{n-1}=1,\qquad d_n=(\lambda-a)^n}.
\]

\newpage
\section{Problem 3}
Find out the Smith canonical form of the $\lambda$-matrix
\[
A(\lambda)=\begin{pmatrix}
\lambda&0&\cdots&0&a_0\\
-1&\lambda&\ddots&\vdots&a_1\\
0&-1&\ddots&0&\vdots\\
\vdots&\ddots&\ddots&\lambda&a_{n-2}\\
0&\cdots&0&-1&\lambda+a_{n-1}
\end{pmatrix}.
\]
\Solution
For $1\le k<n$, the minor on rows $2,\ldots,k+1$ and columns $1,\ldots,k$ equals $(-1)^k$, hence $D_k=1$.
Expanding along the last column gives
\[
D_n=\det A=\lambda^n+a_{n-1}\lambda^{n-1}+\cdots+a_1\lambda+a_0=:p(\lambda).
\]
Consequently,
\[
\boxed{S_A(\lambda)=\operatorname{diag}\bigl(\underbrace{1,\ldots,1}_{n-1},p(\lambda)\bigr)}.
\]

\section{Problem 4}
With respect to a matrix $A\in\F^{n\times n}$ and an indeterminate $\lambda$, $\lambda E-A$ is called the eigenmatrix, where $E$ is the identity matrix. Find out the Smith canonical form of the eigenmatrix of the following matrix
\[
A=\begin{pmatrix}-7&-12&-6\\3&5&3\\3&6&2\end{pmatrix}.
\]
\Solution
If $\operatorname{char}\F\ne3$, including the real and complex fields, put
\[
B(\lambda)=\lambda I-A=\begin{pmatrix}
\lambda+7&12&6\\-3&\lambda-5&-3\\-3&-6&\lambda-2
\end{pmatrix}.
\]
The nonzero constant entry $-3$ gives $D_1=1$. Since $B(-1)$ has rank one, every second-order minor is divisible by $\lambda+1$, while
\[
\begin{vmatrix}\lambda+7&6\\-3&-3\end{vmatrix}=-3(\lambda+1).
\]
Thus $D_2=\lambda+1$. Also $D_3=\det B=(\lambda+1)^2(\lambda-2)$, so
\[
\boxed{S_B(\lambda)=\operatorname{diag}\bigl(1,\lambda+1,(\lambda+1)(\lambda-2)\bigr)}.
\]
If $\operatorname{char}\F=3$, then $A=-I$ and the Smith form is $\boxed{(\lambda+1)I_3}$.
\end{document}
